% GENERATED FILE. Edit canonical lessons, then run npm run generate:books.
% canonical-source-hash: fnv1a64:cfba11717697b036
% canonical-lessons: KA-C466-upanagara, KA-C466-vasisu, KA-C466-purasabhe, KA-C466-bididipa, KA-C466-carandi

\chapter{Living in the Suburbs}
\label{ch:ka-living-in-the-suburbs}
\hlchaptermodality{466}

\begin{chapteropening}
\textbf{By the end of this chapter:} \emph{I can say a suburb, to reside, a town council, a streetlight and a drain.}

Complete the last lesson of chapter 466: I can say a suburb, to reside, a town council, a streetlight and a drain.
\end{chapteropening}

\section[\texorpdfstring{\kn{ಉಪನಗರ} (upanagara)}{ಉಪನಗರ (upanagara)}]{\kn{ಉಪನಗರ} (upanagara) — a suburb}

\label{lesson:KA-C466-upanagara}



\begin{quote}
Before the new one: say the Kannada for a place, then the Kannada for a centre.
\end{quote}

\subsection*{You'll want to know: \kn{ಉಪನಗರ}}
\textbf{\kn{ಉಪನಗರ}} — \emph{upanagara} — \textquotedblleft{}a suburb\textquotedblright{}.

\begin{quote}
\kn{ನಾವು} \kn{ನಗರದ} \kn{ಹೊರಗಿನ} \kn{ಒಂದು} \kn{ಉಪನಗರದಲ್ಲಿ} \kn{ಮನೆ} \kn{ಹುಡುಕಿದೆವು}.
\end{quote}

\begin{quote}
\emph{nāvu nagarada horagina ondu upanagaradalli mane huḍukidevu.}
\end{quote}

\begin{quote}
We looked for a house in a suburb outside the city.
\end{quote}

\kn{ಉಪ}-, near or secondary, plus \kn{ನಗರ}, a city.

\begin{cousinweb}
\textbf{\kn{ಉಪನಗರ}} (\emph{upanagara}) is borrowed from Sanskrit \emph{upanagara}, \textquotedblleft{}a suburb\textquotedblright{} (Monier-Williams), from \emph{upa-}, \textquotedblleft{}near\textquotedblright{}, and \emph{nagara}, \textquotedblleft{}a city\textquotedblright{}.
\end{cousinweb}

\begin{grammarlens}[title={What you've built}]
One more word for getting around.
\end{grammarlens}

\subsection*{Your turn}
\begin{itemize}
\raggedright
  \item \emph{Say it:} \emph{upanagara}
  \item \emph{Say it:} \emph{upanagara}, once more
  \item \emph{Say it:} \emph{nāvu nagarada horagina ondu upanagaradalli mane huḍukidevu.} — We looked for a house in a suburb outside the city.
  \item \emph{Recall:} say the Kannada for a place, then the Kannada for a centre, then say \emph{upanagara} again
\end{itemize}

\begin{tcolorbox}[breakable,colback=teal!4,colframe=teal!35!black,title={Before you move on}]
What does \kn{ಉಪನಗರ} mean? (\textquotedblleft{}A suburb\textquotedblright{}.) Say it once more.
\end{tcolorbox}
\section[\texorpdfstring{\kn{ವಾಸಿಸು} (vāsisu)}{ವಾಸಿಸು (vāsisu)}]{\kn{ವಾಸಿಸು} (vāsisu) — to reside; to live}

\label{lesson:KA-C466-vasisu}



\begin{quote}
Before the new one: say the Kannada for a centre, then the Kannada for a suburb.
\end{quote}

\subsection*{You'll want to know: \kn{ವಾಸಿಸು}}
\textbf{\kn{ವಾಸಿಸು}} — \emph{vāsisu} — \textquotedblleft{}to reside; to live\textquotedblright{}.

\begin{quote}
\kn{ನನ್ನ} \kn{ಅಕ್ಕ} \kn{ಉಪನಗರದಲ್ಲಿ} \kn{ವಾಸಿಸುತ್ತಾಳೆ}.
\end{quote}

\begin{quote}
\emph{nanna akka upanagaradalli vāsisuttāḷe.}
\end{quote}

\begin{quote}
My elder sister lives in a suburb.
\end{quote}

Past (I …): \kn{ವಾಸಿಸಿದೆ}, \emph{vāsiside}.

A formal word for \kn{ಇರು}; \kn{ವಾಸ} is a stay, a residence.

\begin{cousinweb}
\textbf{\kn{ವಾಸಿಸು}} (\emph{vāsisu}) is built from Sanskrit \emph{vāsa}, \textquotedblleft{}an abode, staying\textquotedblright{} (CDIAL 11591), with the Kannada verb ending -\emph{isu}.
\end{cousinweb}

\begin{grammarlens}[title={What you've built}]
One more word for getting around.
\end{grammarlens}

\subsection*{Your turn}
\begin{itemize}
\raggedright
  \item \emph{Say it:} \emph{vāsisu}
  \item \emph{Say it:} \emph{vāsisu}, once more
  \item \emph{Say it:} \emph{nanna akka upanagaradalli vāsisuttāḷe.} — My elder sister lives in a suburb.
  \item \emph{Recall:} say the Kannada for a centre, then the Kannada for a suburb, then say \emph{vāsisu} again
\end{itemize}

\begin{tcolorbox}[breakable,colback=teal!4,colframe=teal!35!black,title={Before you move on}]
What does \kn{ವಾಸಿಸು} mean? (\textquotedblleft{}To reside; to live\textquotedblright{}.) Say it once more.
\end{tcolorbox}
\section[\texorpdfstring{\kn{ಪುರಸಭೆ} (purasabhe)}{ಪುರಸಭೆ (purasabhe)}]{\kn{ಪುರಸಭೆ} (purasabhe) — a town council; a municipality}

\label{lesson:KA-C466-purasabhe}



\begin{quote}
Before the new one: say the Kannada for a suburb, then the Kannada for to reside.
\end{quote}

\subsection*{You'll want to know: \kn{ಪುರಸಭೆ}}
\textbf{\kn{ಪುರಸಭೆ}} — \emph{purasabhe} — \textquotedblleft{}a town council; a municipality\textquotedblright{}.

\begin{quote}
\kn{ಪುರಸಭೆ} \kn{ಬೀದಿಗಳಿಂದ} \kn{ಕಸ} \kn{ಸಂಗ್ರಹಿಸುತ್ತದೆ}.
\end{quote}

\begin{quote}
\emph{purasabhe bīdigaḷinda kasa saṅgrahisuttade.}
\end{quote}

\begin{quote}
The town council collects rubbish from the streets.
\end{quote}

\kn{ಪುರ}, a town, plus \kn{ಸಭೆ}, a meeting or council.

\begin{cousinweb}
\textbf{\kn{ಪುರಸಭೆ}} (\emph{purasabhe}) is built from two Sanskrit words: \emph{pura}, \textquotedblleft{}a fortress, a town\textquotedblright{} (CDIAL 8278), and \emph{sabhā}, \textquotedblleft{}an assembly, a council\textquotedblright{} (CDIAL 13167).
\end{cousinweb}

\begin{grammarlens}[title={What you've built}]
One more word for getting around.
\end{grammarlens}

\subsection*{Your turn}
\begin{itemize}
\raggedright
  \item \emph{Say it:} \emph{purasabhe}
  \item \emph{Say it:} \emph{purasabhe}, once more
  \item \emph{Say it:} \emph{purasabhe bīdigaḷinda kasa saṅgrahisuttade.} — The town council collects rubbish from the streets.
  \item \emph{Recall:} say the Kannada for a suburb, then the Kannada for to reside, then say \emph{purasabhe} again
\end{itemize}

\begin{tcolorbox}[breakable,colback=teal!4,colframe=teal!35!black,title={Before you move on}]
What does \kn{ಪುರಸಭೆ} mean? (\textquotedblleft{}A town council; a municipality\textquotedblright{}.) Say it once more.
\end{tcolorbox}
\section[\texorpdfstring{\kn{ಬೀದಿದೀಪ} (bīdidīpa)}{ಬೀದಿದೀಪ (bīdidīpa)}]{\kn{ಬೀದಿದೀಪ} (bīdidīpa) — a streetlight}

\label{lesson:KA-C466-bididipa}



\begin{quote}
Before the new one: say the Kannada for to reside, then the Kannada for a town council.
\end{quote}

\subsection*{You'll want to know: \kn{ಬೀದಿದೀಪ}}
\textbf{\kn{ಬೀದಿದೀಪ}} — \emph{bīdidīpa} — \textquotedblleft{}a streetlight\textquotedblright{}.

\begin{quote}
\kn{ನಮ್ಮ} \kn{ಬೀದಿಯ} \kn{ಬೀದಿದೀಪ} \kn{ಒಂದು} \kn{ವಾರದಿಂದ} \kn{ಉರಿಯುತ್ತಿಲ್ಲ}.
\end{quote}

\begin{quote}
\emph{namma bīdiya bīdidīpa ondu vāradinda uriyuttilla.}
\end{quote}

\begin{quote}
The streetlight on our street has not been on for a week.
\end{quote}

\kn{ಬೀದಿ}, a street, plus \kn{ದೀಪ}, a lamp.

\begin{cousinweb}
\textbf{\kn{ಬೀದಿದೀಪ}} (\emph{bīdidīpa}) is a Kannada compound built from \kn{ಬೀದಿ}, \textquotedblleft{}a street\textquotedblright{}, from Sanskrit \emph{vīthi}, \textquotedblleft{}a row, a road\textquotedblright{} (Monier-Williams), and \kn{ದೀಪ}, \textquotedblleft{}a lamp\textquotedblright{} (CDIAL 6348).
\end{cousinweb}

\begin{grammarlens}[title={What you've built}]
One more word for getting around.
\end{grammarlens}

\subsection*{Your turn}
\begin{itemize}
\raggedright
  \item \emph{Say it:} \emph{bīdidīpa}
  \item \emph{Say it:} \emph{bīdidīpa}, once more
  \item \emph{Say it:} \emph{namma bīdiya bīdidīpa ondu vāradinda uriyuttilla.} — The streetlight on our street has not been on for a week.
  \item \emph{Recall:} say the Kannada for to reside, then the Kannada for a town council, then say \emph{bīdidīpa} again
\end{itemize}

\begin{tcolorbox}[breakable,colback=teal!4,colframe=teal!35!black,title={Before you move on}]
What does \kn{ಬೀದಿದೀಪ} mean? (\textquotedblleft{}A streetlight\textquotedblright{}.) Say it once more.
\end{tcolorbox}
\section[\texorpdfstring{\kn{ಚರಂಡಿ} (caraṇḍi)}{ಚರಂಡಿ (caraṇḍi)}]{\kn{ಚರಂಡಿ} (caraṇḍi) — a drain}

\label{lesson:KA-C466-carandi}



\begin{quote}
Before the new one: say the Kannada for a town council, then the Kannada for a streetlight.
\end{quote}

\subsection*{You'll want to know: \kn{ಚರಂಡಿ}}
\textbf{\kn{ಚರಂಡಿ}} — \emph{caraṇḍi} — \textquotedblleft{}a drain\textquotedblright{}.

\begin{quote}
\kn{ಮಳೆಯಲ್ಲಿ} \kn{ಚರಂಡಿ} \kn{ತುಂಬಿ} \kn{ನೀರು} \kn{ರಸ್ತೆಗೆ} \kn{ಬಂದಿತು}.
\end{quote}

\begin{quote}
\emph{maḷeyalli caraṇḍi tumbi nīru rastege banditu.}
\end{quote}

\begin{quote}
In the rain the drain filled up and water came onto the road.
\end{quote}

The open drain beside a street; \kn{ಒಳಚರಂಡಿ} is an underground sewer.

\begin{cousinweb}
\textbf{\kn{ಚರಂಡಿ}} (\emph{caraṇḍi}) is a Kannada word for a drain whose origin is not settled.
\end{cousinweb}

\begin{grammarlens}[title={What you've built}]
One more word for getting around.
\end{grammarlens}

\subsection*{Your turn}
\begin{itemize}
\raggedright
  \item \emph{Say it:} \emph{caraṇḍi}
  \item \emph{Say it:} \emph{caraṇḍi}, once more
  \item \emph{Say it:} \emph{maḷeyalli caraṇḍi tumbi nīru rastege banditu.} — In the rain the drain filled up and water came onto the road.
  \item \emph{Recall:} say the Kannada for a town council, then the Kannada for a streetlight, then say \emph{caraṇḍi} again
\end{itemize}

\begin{tcolorbox}[breakable,colback=teal!4,colframe=teal!35!black,title={Before you move on}]
What does \kn{ಚರಂಡಿ} mean? (\textquotedblleft{}A drain\textquotedblright{}.) Say it once more.
\end{tcolorbox}
